% Generated by explore_fda/make_paper_outputs.R -- regenerate rather than edit
\begin{table}[!htbp]
\centering
\begin{threeparttable}
\caption{Group differences in athlete-mean force curves}
\label{tab:spm}
\footnotesize
\begin{tabular*}{\linewidth}{@{\extracolsep{\fill}}l l c S[table-format=2.0] S[table-format=1.2] l S[table-format=1.2] l@{}}
\toprule
 & & & & \multicolumn{2}{c}{Difference (BW)} & \multicolumn{2}{c}{Hedges' \textit{g}} \\
\cmidrule(lr){5-6}\cmidrule(l){7-8}
Phase & {Region (\%)} & {\textit{p}} & {Peak (\%)} & {Value} & {95\% CI} & {Value} & {95\% CI} \\
\midrule
\multicolumn{8}{@{}l}{\textit{Panel A. Top vs bottom tertile of jump height (71 vs 71 athletes)}} \\
\quad Take-off & {81--99} & {$<$.001} & 93 & 0.55 & {[0.48, 0.61]} & 2.63 & {[2.31, 2.94]} \\
\quad Landing & {2--56} & .003 & 15 & 0.75 & {[0.45, 1.04]} & 0.83 & {[0.50, 1.15]} \\
\multicolumn{8}{@{}l}{\hspace{0.5em}\textit{Adjusted for age and body mass}$^{a}$} \\
\quad Take-off & {12--32} & .021 & {--} & {--} & {--} & {--} & {--} \\
\quad  & {66--99} & {$<$.001} & {--} & {--} & {--} & {--} & {--} \\
\quad Landing & {19--40} & .043 & {--} & {--} & {--} & {--} & {--} \\
\addlinespace
\multicolumn{8}{@{}l}{\textit{Panel B. Aged 18 or over vs under 18 (48 vs 164 athletes)}} \\
\quad Take-off & {87--99} & {$<$.001} & 96 & 0.47 & {[0.36, 0.56]} & 1.60 & {[1.25, 1.92]} \\
\quad Landing & {none} & .069$^{c}$ & {--} & {--} & {--} & {--} & {--} \\
\multicolumn{8}{@{}l}{\hspace{0.5em}\textit{Adjusted for body mass and jump height}$^{a}$} \\
\quad Take-off & {none} & .123$^{c}$ & {--} & {--} & {--} & {--} & {--} \\
\quad Landing & {none} & .226$^{c}$ & {--} & {--} & {--} & {--} & {--} \\
\bottomrule
\end{tabular*}
\begin{tablenotes}[flushleft]\footnotesize
\item \textit{Note.} One mean curve per athlete. Region = cluster of adjacent nodes (\% of the phase) where the groups differed; \textit{p} from 5000 permutations of the largest cluster mass, controlling the family-wise error rate over the phase$^{b}$. Peak = node of the largest |\textit{t}| in the region, where the difference (first group minus second; BW = multiples of body weight) and Hedges' \textit{g} are given with 95\% confidence intervals (CIs) from 2000 bootstrap resamples of athletes within groups. Bonferroni threshold for the four primary tests: \textit{p} < .0125.
\item $^{a}$Freedman--Lane permutation of residuals. $^{b}$Nodes with $|t|$ above the two-sided 5\% threshold form clusters. $^{c}$No significant region; smallest cluster \textit{p} shown.
\end{tablenotes}
\end{threeparttable}
\end{table}
